%% Multi-level and graph-structured generalisation of the hierarchical
%% composition rule (gauss_hierarchical / gauss_jacobian).
%% Derivation only -- nothing here is implemented in src/compass.
\documentclass[11pt]{article}
\usepackage[a4paper,margin=2.4cm]{geometry}
\usepackage{amsmath,amssymb,amsthm,mathtools,bm}
\usepackage{microtype,booktabs,xcolor}
\usepackage[T1]{fontenc}
\usepackage{lmodern}
\usepackage[colorlinks=true,linkcolor=blue!50!black,citecolor=blue!50!black]{hyperref}
\usepackage[capitalize,noabbrev]{cleveref}

\newtheorem{theorem}{Theorem}
\newtheorem{proposition}[theorem]{Proposition}
\newtheorem{lemma}[theorem]{Lemma}
\newtheorem{corollary}[theorem]{Corollary}
\theoremstyle{definition}
\newtheorem{definition}{Definition}
\newtheorem{assumption}{Assumption}
\crefname{assumption}{Assumption}{Assumptions}
\newtheorem{example}{Example}
\theoremstyle{remark}
\newtheorem{remark}{Remark}

\newcommand{\R}{\mathbb{R}}
\newcommand{\E}{\mathbb{E}}
\newcommand{\N}{\mathcal{N}}
\newcommand{\I}{\mathbf{I}}
\newcommand{\T}{\top}
\newcommand{\dd}{\,\mathrm{d}}
\newcommand{\lam}{\lambda}
\newcommand{\X}{\mathbf{x}}
\newcommand{\An}{\operatorname{An}}
\newcommand{\pa}{\operatorname{pa}}
\newcommand{\ch}{\operatorname{ch}}
\newcommand{\sub}{\operatorname{sub}}
\newcommand{\Rset}{\mathcal{R}}
\newcommand{\Fset}{\mathcal{F}}
\newcommand{\arrowm}{\texttt{gauss\_hierarchical}}
\newcommand{\quiver}{\texttt{gauss\_jacobian}}

\title{Composing scores beyond two levels:\\
three-level, $L$-level and graph-structured hierarchies}
\author{COMPASS working note (derivation only, no code changes)}
\date{12 October 2026}

\begin{document}
\maketitle

\begin{abstract}
The hierarchical composition rule behind \arrowm{} and \quiver{} composes a
single-observation score network over a two-level model (shared $g$, one local
$\ell_j$ per observation). We show that the rule generalises without new
approximations. The same single Gaussian-backward-kernel approximation gives
(i)~an explicit three-level rule, (ii)~an $L$-level rule for any rooted tree,
obtained by applying the two-level rule recursively (``a group is a big
observation''), and (iii)~a rule for any hierarchical model whose latent
structure is a directed acyclic graph (DAG) and whose observations are sinks. In
all three cases the tall-data backward kernel factorises \emph{exactly} into
single-observation kernels and prior-marginal kernels, weighted by integer
\emph{counting numbers}. The Gaussian step then reduces the composed score to a
sparse linear system $\Lambda u=r$ assembled only from the network's own
(conditional) scores and precisions. On trees, block elimination from the leaves
up solves the system in time linear in the number of nodes, with no fill-in.
On general DAGs the system is still exact, but its cost is set by the treewidth
of a derived graph. The rule needs one new ingredient from three levels on: the
diffused prior score and backward precision on every intermediate
\emph{path} (e.g.\ $p_t(g,h_k)$), not only on $g$. A NumPy check on
linear-Gaussian trees (3 and 5 levels), a crossed-effects design and a non-tree
DAG reproduces the exact diffused tall-data score to machine precision
($\le 7\times10^{-15}$ relative error) at every noise level. Dropping the
intermediate prior corrections gives 1--67\% error.
\end{abstract}

\tableofcontents

%=====================================================================
\section{Recap: the two-level rule}\label{sec:recap}
%=====================================================================

The current model is
$p(g,\ell_{1:n},\X_{1:n})=p(g)\prod_{j}p(\ell_j\mid g)\,p(x_j\mid g,\ell_j)$.
The network provides, for each observation, the joint score
$s_j=(s_j^g,s_j^\ell)$ of $p_t(\vartheta_j\mid x_j)$ at the noisy
$y_j=(g_t,\ell_{j,t})$, with $\vartheta_j=(g,\ell_j)$. Let $\varphi_\lam$ be the
isotropic Gaussian kernel at noise scale $\lam=\lam_t$ (VE convention, as in the
paper). The derivation (paper, \S3) has four steps:
\begin{enumerate}
\item \textbf{Exact factorisation:}
 $p(\theta\mid\X)\propto p(g)^{1-n}\prod_jp(\vartheta_j\mid x_j)$. Because the
 forward kernel factorises over coordinates, the same holds for the backward
 kernels:
 $\pi_t(\theta_0\mid y,\X)\propto\pi_t^0(g_0\mid g_t)^{1-n}\prod_j\pi_t(\vartheta_{j,0}\mid y_j,x_j)$.
\item \textbf{Tweedie:} the composed score is $u=(\E[\theta_0\mid y,\X]-y)/\lam^2$.
\item \textbf{Gaussian factors:} replacing each factor by a Gaussian with the
 correct (Tweedie) mean and precision $P_j$ (resp.\ $P_0$) gives $\Lambda u=r$
 with
 $\Lambda=(1-n)E_g^\T P_0E_g+\sum_jE_j^\T P_jE_j$ and
 $r=(1-n)E_g^\T P_0s_0+\sum_jE_j^\T P_js_j$. $\Lambda$ is an arrowhead matrix.
\item \textbf{Elimination:} with $P_j=\begin{psmallmatrix}A_j&B_j\\B_j^\T&C_j\end{psmallmatrix}$,
 $S_j=A_j-B_jC_j^{-1}B_j^\T$ and $R_j=-C_j^{-1}B_j^\T$,
 \begin{equation}\label{eq:two-level}
 u^g=\Big[\sum_jS_j+(1-n)P_0\Big]^{-1}\Big[\sum_jS_js_j^g+(1-n)P_0s_0\Big],\qquad
 u^{\ell_j}=s_j^\ell+R_j(u^g-s_j^g).
 \end{equation}
\end{enumerate}
The rest of this note reuses steps 2 and 3 unchanged. All the work goes into
(a)~finding the right exact factorisation (step~1) for deeper or graph-structured
models, and (b)~solving the resulting larger sparse system (step~4).

\paragraph{Two observations that make everything below work.}
First, steps 2--3 never used the arrowhead structure. They hold for \emph{any}
finite product of Gaussian factors with real (possibly negative) exponents,
provided every coordinate of $y$ receives total noise exponent one (\cref{prop:linear-system}).
Second, step~4 is a single instance of a generic operation, the elimination of a
block that only one (merged) factor touches (\cref{lem:marginalise}). Applying it
repeatedly gives the multi-level rules.

%=====================================================================
\section{General setting and the exact factorisation}\label{sec:general}
%=====================================================================

\subsection{Model class}

\begin{assumption}[Hierarchical model on a DAG]\label{ass:dag}
The latent variables are blocks $\theta_v\in\R^{d_v}$ indexed by the nodes
$v\in V$ of a DAG $G$ with parent sets $\pa(v)$, and
\[
p(\theta,\X)=\prod_{v\in V}p(\theta_v\mid\theta_{\pa(v)})\;\prod_{o\in O}p(x_o\mid\theta_{\pa(o)}),
\]
where every observation $o\in O$ is a sink: $\pa(o)\subseteq V$, and no latent
or other observation has $o$ as a parent. Every latent node has at least one
observed descendant (latent nodes without one are barren; integrate them out
first, which leaves the model in the same class).
\end{assumption}
The two-level model is the DAG $g\to\ell_j\to x_j$ (with also $g\to x_j$). The
three-level model is $g\to h_k\to\ell_{kj}\to x_{kj}$. Crossed random effects
$g\to a_i$, $g\to b_j$, $(a_i,b_j)\to x_{ij}$ are a DAG that is not a tree.

\begin{definition}[Observation regions]
For $o\in O$, let $R_o:=\An(\pa(o))$ be the ancestral closure of its parents:
all latent nodes from which there is a directed path into $o$, including
$\pa(o)$. Write $\theta_R=(\theta_v)_{v\in R}$ and $d_R=\sum_{v\in R}d_v$.
\end{definition}

\begin{lemma}[Ancestral marginals are products of CPDs]\label{lem:ancestral}
For any ancestral set $R\subseteq V$ (closed under taking parents),
$p(\theta_R)=\prod_{v\in R}p(\theta_v\mid\theta_{\pa(v)})$.
\end{lemma}
\begin{proof}
Integrate out the nodes of $V\setminus R$ in reverse topological order. Each one
integrated out is a sink of the current graph, since its children are outside
$R$ (otherwise $R$ would not be ancestral) and have already been removed. Its
CPD therefore integrates to one, and no remaining factor depends on it.
\end{proof}

The training unit of the network is consequently \emph{one observation together
with its ancestral region}: draw $\theta_{R_o}$ by ancestral sampling, then
$x_o$. The network learns $p(\theta_{R_o}\mid x_o)$ under the correct prior
marginal $p(\theta_{R_o})$, and no simulation ever involves more than one
observation. In the two-level model $R_j=\{g,\ell_j\}$, which is the current
single-group training unit. In the three-level model $R_{kj}=\{g,h_k,\ell_{kj}\}$,
the root-to-leaf path.

\subsection{Counting numbers}

\begin{definition}[Region family and counting numbers]\label{def:counting}
Let $\Rset$ be the closure of $\{R_o:o\in O\}$ under non-empty intersection.
Intersections of ancestral sets are ancestral, so every $R\in\Rset$ is ancestral.
Define the (Möbius) counting numbers top-down,
\begin{equation}\label{eq:moebius}
c(R):=1-\sum_{R'\in\Rset,\;R'\supsetneq R}c(R'),
\end{equation}
let $m(R):=|\{o:R_o=R\}|$ be the number of observations whose region is $R$, and
define the \emph{prior exponent} $\kappa(R):=c(R)-m(R)$.
\end{definition}

\begin{lemma}[Every node is covered exactly once]\label{lem:cover}
For every latent node $v$,
\begin{equation}\label{eq:cover}
\underbrace{|\{o\in O:v\in R_o\}|}_{\text{observation factors}}+\sum_{R\in\Rset,\,R\ni v}\kappa(R)=\sum_{R\in\Rset,\,R\ni v}c(R)=1.
\end{equation}
\end{lemma}
\begin{proof}
The first equality is the definition of $\kappa$. For the second, let
$S_v:=\bigcap\{R\in\Rset:v\in R\}$. The set is non-empty by
\cref{ass:dag}, and the intersection contains $v$, so $S_v\in\Rset$. A region
contains $v$ if and only if it contains $S_v$. Hence
$\sum_{R\ni v}c(R)=c(S_v)+\sum_{R\supsetneq S_v}c(R)=1$ by \eqref{eq:moebius}.
\end{proof}

\begin{theorem}[Exact tall-data factorisation on a DAG]\label{thm:factorisation}
Under \cref{ass:dag},
\begin{equation}\label{eq:factorisation}
p(\theta\mid\X)\;\propto\;\prod_{o\in O}p(\theta_{R_o}\mid x_o)\;\prod_{R\in\Rset}p(\theta_R)^{\kappa(R)},
\end{equation}
and for every $t$ the backward kernel of the diffused tall posterior factorises
exactly as
\begin{equation}\label{eq:backward-factorisation}
\pi_t(\theta_0\mid y,\X)\;\propto\;\prod_{o\in O}\pi_t(\theta_{R_o,0}\mid y_{R_o},x_o)\;\prod_{R\in\Rset}\pi^0_t(\theta_{R,0}\mid y_R)^{\kappa(R)},
\end{equation}
where $\pi^0_t(\theta_{R,0}\mid y_R)\propto\varphi_\lam(y_R-\theta_{R,0})\,p(\theta_{R,0})$
is the backward kernel of the prior marginal on $R$.
\end{theorem}
\begin{proof}
\emph{Clean posterior.} By Bayes and \cref{lem:ancestral},
$p(\theta_{R_o}\mid x_o)\propto p(x_o\mid\theta_{\pa(o)})\prod_{v\in R_o}p(\theta_v\mid\theta_{\pa(v)})$,
and $p(\theta_R)=\prod_{v\in R}p(\theta_v\mid\theta_{\pa(v)})$. In the right-hand side of
\eqref{eq:factorisation} each likelihood appears once. The CPD of node $v$
appears with total exponent equal to the left-hand side of \eqref{eq:cover},
which is $1$. The right-hand side is therefore proportional to the joint
$p(\theta,\X)$, hence to $p(\theta\mid\X)$.

\emph{Backward kernel.} Since $\varphi_\lam(y-\theta_0)=\prod_v\varphi_\lam(y_v-\theta_{v,0})$
and, by \eqref{eq:cover}, the exponents of the copies handed to the factors sum
to one for each $v$,
\[
\varphi_\lam(y-\theta_0)=\prod_{o}\varphi_\lam(y_{R_o}-\theta_{R_o,0})\prod_{R}\varphi_\lam(y_R-\theta_{R,0})^{\kappa(R)}.
\]
Multiply this by \eqref{eq:factorisation} and pair each noise factor with its
density factor. Each pair is a backward kernel times a normaliser
($p_t(y_{R_o}\mid x_o)$ or $p_t(y_R)^{\kappa(R)}$) that depends on $y$ but not
on $\theta_0$.
\end{proof}

\begin{remark}
The counting condition \eqref{eq:cover} has the same form as the counting-number
condition of region-based (Kikuchi) approximations in graphical models
(Yedidia, Freeman \& Weiss, 2005). Here, however, it is \emph{exact}, because the
regions are ancestral (\cref{lem:ancestral}). The Möbius choice
\eqref{eq:moebius} is canonical but not unique. Any integer $\kappa$ satisfying
\eqref{eq:cover} gives an exact factorisation; different solutions can differ in
numerical conditioning once the factors are approximated.
\end{remark}

\subsection{From the factorisation to a linear system}

\begin{proposition}[Composed score as a sparse linear system]\label{prop:linear-system}
Index all factors in \eqref{eq:backward-factorisation} by $f\in\Fset$. Each has
weight $w_f$ ($1$ for an observation factor, $\kappa(R)$ for a prior factor),
scope $R_f$ with selector $E_f$ ($E_f\theta=\theta_{R_f}$), score
$s_f=\nabla\log p_t(y_{R_f}\mid\cdot)$ and a Gaussian surrogate backward
precision $P_f\succ0$. If
\begin{equation}\label{eq:Lambda}
\Lambda:=\sum_{f\in\Fset}w_f\,E_f^\T P_fE_f\;\succ0,\qquad r:=\sum_{f\in\Fset}w_f\,E_f^\T P_fs_f,
\end{equation}
then the Gaussian surrogate of \eqref{eq:backward-factorisation} gives the
composed score $u=\nabla_y\log p_t(y\mid\X)\approx\Lambda^{-1}r$. Given exact
factor scores and precisions, this is exact whenever all factors are Gaussian, in
particular for every linear-Gaussian DAG.
\end{proposition}
\begin{proof}
Verbatim the paper's Step~3 with $(1-n)$ and $1$ replaced by $w_f$. With
$m_f=y_{R_f}+\lam^2s_f$ (Tweedie), the log surrogate is
$-\tfrac12\theta_0^\T\Lambda\theta_0+\theta_0^\T b$ with
$b=\sum_fw_fE_f^\T P_fm_f$. Its mean is $\Lambda^{-1}b$, and
$u=(\Lambda^{-1}b-y)/\lam^2=\Lambda^{-1}(b-\Lambda y)/\lam^2$. Since
$E_fy=y_{R_f}$, $b-\Lambda y=\sum_fw_fE_f^\T P_f(m_f-y_{R_f})=\lam^2r$.
\end{proof}

The \emph{sparsity graph} of $\Lambda$ connects $v$ and $v'$ only if they share a
region. It is therefore the union of the cliques $\{R_o\}$, which we call the
\emph{ancestral observation graph}. All structural results below concern how to
solve $\Lambda u=r$ efficiently on this graph. The \emph{values} of $u$ do not
depend on how it is solved.

\subsection{The elimination lemma}

\begin{lemma}[Merge and marginalise]\label{lem:marginalise}
\textbf{(Merge.)} Factors with the same scope can be summed. Weighted factors
$(w_i,P_i,s_i)$ on one scope are equivalent in $\Lambda,r$ to one factor
$(Q,q)$ with $Q=\sum_iw_iP_i$ and $q=\sum_iw_iP_is_i$. If $Q\succ0$, write
$q=Q\sigma$ with the \emph{merged score} $\sigma:=Q^{-1}q$.

\textbf{(Marginalise.)} Let a block $c$ of the unknowns appear in exactly one
(merged) factor $(Q,Q\sigma)$, with scope $(a,c)$ and
$Q=\begin{psmallmatrix}Q_{aa}&Q_{ac}\\Q_{ca}&Q_{cc}\end{psmallmatrix}$,
$Q_{cc}\succ0$. Replacing this factor by the factor $(S,\,S\sigma^a)$ on scope
$a$ alone, with
\begin{equation}\label{eq:schur}
S=Q_{aa}-Q_{ac}Q_{cc}^{-1}Q_{ca}=\big([Q^{-1}]_{aa}\big)^{-1},
\end{equation}
and back-substituting
\begin{equation}\label{eq:backsub}
u^c=\sigma^c+K\,(u^a-\sigma^a),\qquad K:=-Q_{cc}^{-1}Q_{ca},
\end{equation}
leaves the solution $u$ of $\Lambda u=r$ unchanged.
\end{lemma}
\begin{proof}
Merging holds because $\Lambda$ and $r$ are sums over factors. For
marginalising, the $c$-rows of $\Lambda u=r$ read
$Q_{ca}u^a+Q_{cc}u^c=Q_{ca}\sigma^a+Q_{cc}\sigma^c$, which gives
\eqref{eq:backsub}. The factor's contribution to the $a$-rows is
$Q_{aa}u^a+Q_{ac}u^c$ on the left and $Q_{aa}\sigma^a+Q_{ac}\sigma^c$ on the
right. Substituting $u^c$ cancels every $\sigma^c$ term and leaves
$Su^a=S\sigma^a$. Nothing else in the system involves $c$.
\end{proof}

This is exactly the paper's Step~4. In the two-level rule, $c=\ell_j$,
$a=g$, the factor is $(P_j,s_j)$, and the merged score is the network score
$s_j$ itself, so $S=S_j$, $K=R_j$ and the message is $(S_j,s_j^g)$. The
``cancellation of the local scores'' is the statement that the message keeps
$\sigma^a$, the \emph{conditional} score of the factor at the current $c$. The
marginal score of the scope $a$ is never required. The lemma shows this
cancellation happens at \emph{every} elimination, at any depth and in any
graph.

%=====================================================================
\section{Three levels}\label{sec:three}
%=====================================================================

\subsection{Model and exact factorisation}

\[
p(\theta,\X)=p(g)\prod_{k=1}^{K}\Big[p(h_k\mid g)\prod_{j=1}^{n_k}p(\ell_{kj}\mid g,h_k)\,p(x_{kj}\mid g,h_k,\ell_{kj})\Big],
\]
with $g\in\R^{d_g}$ (global), $h_k\in\R^{d_h}$ (group), $\ell_{kj}\in\R^{d_\ell}$
(local) and $N=\sum_kn_k$ observations. Example: a population mean $g$, school
effects $h_k$, pupil effects $\ell_{kj}$, and pupil data $x_{kj}$. The network is
trained on paths $\vartheta_{kj}=(g,h_k,\ell_{kj})\in\R^{D}$, $D=d_g+d_h+d_\ell$,
and returns $s_{kj}=(s^g_{kj},s^h_{kj},s^\ell_{kj})$ at
$y_{kj}=(g_t,h_{k,t},\ell_{kj,t})$.

The regions are $R_{kj}=\{g,h_k,\ell_{kj}\}$. Their intersections are $\{g,h_k\}$
(within a group) and $\{g\}$ (across groups). \Cref{def:counting} gives
$c(R_{kj})=1$, $c(\{g,h_k\})=1-n_k$ and $c(\{g\})=1-\sum_k[n_k+(1-n_k)]=1-K$.
Hence
\begin{equation}\label{eq:three-factorisation}
p(\theta\mid\X)\propto p(g)^{1-K}\prod_{k=1}^{K}p(g,h_k)^{1-n_k}\prod_{k,j}p(\vartheta_{kj}\mid x_{kj}).
\end{equation}
A direct check: $p(g)$ appears with exponent $(1-K)+\sum_k(1-n_k)+\sum_kn_k=1$,
$p(h_k\mid g)$ with exponent $(1-n_k)+n_k=1$, and $p(\ell_{kj}\mid g,h_k)$
once. The same exponents distribute the noise kernel, so by
\cref{thm:factorisation}
\[
\pi_t(\theta_0\mid y,\X)\propto\pi_t^0(g_0\mid g_t)^{1-K}\prod_k\pi_t^0(g_0,h_{k,0}\mid g_t,h_{k,t})^{1-n_k}\prod_{k,j}\pi_t(\vartheta_{kj,0}\mid y_{kj},x_{kj}).
\]

\paragraph{What is new compared with two levels.} The middle factor
$p(g,h_k)^{1-n_k}$ is a correction on the \emph{joint} prior marginal of the
group path. Its diffused score $s_{0,k}=\nabla\log p_t(g_t,h_{k,t})$ and backward
precision $P_{0,k}\in\R^{(d_g+d_h)^2}$ are needed. The conditional prior
$p(h_k\mid g)$ alone is not enough, because diffusion acts on $(g,h_k)$ jointly
and $p_t(g,h_k)\neq p_t(g)\,p_t(h_k\mid g)$ for any useful sense of the latter.
With two levels, the corresponding factor $p(g,\ell_j)^{1-1}$ has exponent zero,
which is why it never showed up.

\subsection{Linear system}

Denote the top prior factor by $(P_0,s_0)$ on $g$, the group prior factors by
$(P_{0,k},s_{0,k})$ on $a_k:=(g,h_k)$, and the observation factors by
$(P_{kj},s_{kj})$ on $\vartheta_{kj}$. Then \eqref{eq:Lambda} reads
\[
\Lambda=(1-K)E_g^\T P_0E_g+\sum_k(1-n_k)E_{a_k}^\T P_{0,k}E_{a_k}+\sum_{k,j}E_{kj}^\T P_{kj}E_{kj},
\]
and $r$ is the same with each $P$ replaced by $Ps$. Ordered as
$(g;\,h_1,\ell_{1\cdot};\,\dots;\,h_K,\ell_{K\cdot})$, $\Lambda$ is a
\emph{nested arrowhead}. The $g$-row couples to everything. Each $h_k$ couples
to $g$ and to its own $\ell_{k\cdot}$. Each $\ell_{kj}$ couples only to $g$ and
$h_k$.

\subsection{Elimination in three passes}

Split each observation precision with respect to $(a_k,\ell_{kj})$ as
$P_{kj}=\begin{psmallmatrix}A_{kj}&B_{kj}\\B_{kj}^\T&C_{kj}\end{psmallmatrix}$, with
$A_{kj}\in\R^{(d_g+d_h)^2}$.

\paragraph{Pass 1: locals.} $\ell_{kj}$ appears only in factor $kj$.
\Cref{lem:marginalise} with $c=\ell_{kj}$, $a=a_k$ gives
\begin{equation}\label{eq:pass1}
S_{kj}=A_{kj}-B_{kj}C_{kj}^{-1}B_{kj}^\T=\big([\Sigma_{t,kj}]_{a a}\big)^{-1},\quad
K_{kj}=-C_{kj}^{-1}B_{kj}^\T,\quad
u^{\ell_{kj}}=s^\ell_{kj}+K_{kj}\big(u^{a_k}-s^{a}_{kj}\big),
\end{equation}
where $s^a_{kj}=(s^g_{kj},s^h_{kj})$. Factor $kj$ becomes the message
$(S_{kj},s^a_{kj})$ on $a_k=(g,h_k)$.

\paragraph{Pass 2: groups.} All remaining factors involving $h_k$ now have scope
$a_k$: the $n_k$ messages and the group prior. Merge them:
\begin{equation}\label{eq:pass2-merge}
Q_k=\sum_{j=1}^{n_k}S_{kj}+(1-n_k)P_{0,k},\qquad
q_k=\sum_{j=1}^{n_k}S_{kj}s^a_{kj}+(1-n_k)P_{0,k}s_{0,k},\qquad
\sigma_k=Q_k^{-1}q_k.
\end{equation}
$\sigma_k$ is exactly the two-level rule \eqref{eq:two-level} applied to group $k$
alone, with ``global'' $a_k=(g,h_k)$, locals $\ell_{k\cdot}$ and prior $p(g,h_k)$.
It is the composed score of group $k$'s tall posterior
$p(g,h_k,\ell_{k\cdot}\mid x_{k\cdot})$ on its $(g,h_k)$ block. Now split
$Q_k=\begin{psmallmatrix}Q_k^{gg}&Q_k^{gh}\\Q_k^{hg}&Q_k^{hh}\end{psmallmatrix}$ and eliminate $h_k$:
\begin{equation}\label{eq:pass2}
T_k=Q_k^{gg}-Q_k^{gh}(Q_k^{hh})^{-1}Q_k^{hg}=\big([Q_k^{-1}]_{gg}\big)^{-1},\quad
G_k=-(Q_k^{hh})^{-1}Q_k^{hg},\quad
u^{h_k}=\sigma^h_k+G_k\big(u^g-\sigma^g_k\big).
\end{equation}
Group $k$ becomes the message $(T_k,\sigma_k^g)$ on $g$.

\paragraph{Pass 3: global.}
\begin{equation}\label{eq:pass3}
u^g=\Big[\sum_{k=1}^{K}T_k+(1-K)P_0\Big]^{-1}\Big[\sum_{k=1}^{K}T_k\,\sigma^g_k+(1-K)P_0\,s_0(g_t)\Big].
\end{equation}

\paragraph{The three-level rule.} Compute $S_{kj},K_{kj}$ (pass 1), then
$Q_k,\sigma_k,T_k,G_k$ (pass 2), then $u^g$ \eqref{eq:pass3}. Back-substitute
$u^{h_k}$ by \eqref{eq:pass2} and $u^{\ell_{kj}}$ by \eqref{eq:pass1}, using
$u^{a_k}=(u^g,u^{h_k})$.

Read top-down, the global row \eqref{eq:pass3} is the two-level rule with the
\emph{groups} as observations. Each group enters with its composed score
$\sigma^g_k$ and its effective precision $T_k$. Each group score is in turn the
two-level rule over its own observations. The three-level rule is therefore the
two-level rule nested in itself. The cost is $N$ inversions of $d_\ell\times d_\ell$
blocks, $K$ of $(d_g+d_h)\times(d_g+d_h)$ blocks, and one $d_g\times d_g$ solve:
linear in $N$.

\begin{remark}[Numerically preferred form]
The merged score $\sigma_k$ needs $Q_k$ to be invertible. To eliminate $h_k$,
only $Q^{hh}_k\succ0$ is needed. Without forming $\sigma_k$, the message
right-hand side is $\tau_k=q_k^g-Q_k^{gh}(Q_k^{hh})^{-1}q_k^h$ (which equals
$T_k\sigma_k^g$ when $Q_k$ is invertible), and the back-substitution is
$u^{h_k}=(Q_k^{hh})^{-1}(q_k^h-Q_k^{hg}u^g)$. Implementations should use this
information form.
\end{remark}

\subsection{Consistency checks}

Each check follows because the rule returns $\Lambda^{-1}r$ and $(\Lambda,r)$
coincide with the reference.
\begin{enumerate}
\item \textbf{One observation per group} ($n_k=1$): the group priors have
 exponent $0$, $Q_k=S_{k1}$ and $\sigma_k=s^a_{k1}$. The rule is the two-level
 rule with global $g$ and local block $(h_k,\ell_{k1})$.
\item \textbf{One group} ($K=1$): the top prior has exponent $0$ and
 $u^g=\sigma^g_1$. The rule is the two-level rule with global $(g,h_1)$ and prior
 $p(g,h_1)$.
\item \textbf{No group parameters} ($d_h=0$): $P_{0,k}=P_0$ and $T_k=Q_k$, so
 \eqref{eq:pass3} becomes
 $\big[\sum_{kj}S_{kj}+(1-N)P_0\big]^{-1}\big[\sum_{kj}S_{kj}s^g_{kj}+(1-N)P_0s_0\big]$,
 which is the two-level rule over all $N$ observations.
\item \textbf{$N=1$:} every prior exponent vanishes, and the rule returns the network score.
\item \textbf{Linear-Gaussian model:} exact at every $t$ (\cref{prop:linear-system}); verified in \cref{sec:numerics}.
\end{enumerate}

%=====================================================================
\section{$L$ levels and arbitrary trees}\label{sec:tree}
%=====================================================================

\subsection{Setting}
Let the latent DAG be a rooted tree: every $v$ has at most one parent, and the
root $\rho$ has none. Observations may attach to any node, provided all parents
of an observation lie on one root path. Write $O(v)$ for the
observations whose deepest parent is $v$ (so $R_o=\pi(v)$), $\ch(v)$ for the
children of $v$ and $\sub(v)$ for the subtree rooted at $v$. Let $\pi(v)$ be the
path from $\rho$ to $v$ (inclusive) and $\alpha(v):=\pi(v)\setminus\{v\}$ the
strict ancestors, so that $\alpha(w)=\pi(v)$ for every $w\in\ch(v)$. This covers
$L$-level nested designs, ragged trees, and observations on intermediate levels.

\begin{proposition}[Counting numbers on a tree]\label{prop:tree-counting}
On a tree, $\Rset=\{\pi(v)\}$, and the prior exponent of $\pi(v)$ is
\begin{equation}\label{eq:tree-kappa}
\kappa(v):=1-|\ch(v)|-|O(v)|.
\end{equation}
\end{proposition}
\begin{proof}
Intersections of root paths are root paths. Given \eqref{eq:tree-kappa}, check
\eqref{eq:cover} by induction from the leaves: let
$N(v):=\sum_{u\in\sub(v)}|O(u)|$ be the number of observations below $v$. Then
\begin{equation}\label{eq:subtree-sum}
N(v)+\sum_{u\in\sub(v)}\kappa(u)=|O(v)|+\kappa(v)+\sum_{w\in\ch(v)}\Big[N(w)+\sum_{u\in\sub(w)}\kappa(u)\Big]=|O(v)|+\kappa(v)+|\ch(v)|=1.
\end{equation}
The left-hand side of \eqref{eq:subtree-sum} is exactly the coverage of $v$, and
also of every ancestor of $v$, contributed by $\sub(v)$. Applied at the root,
it shows every node is covered once. The Möbius numbers satisfy the same
recursion, so they coincide with \eqref{eq:tree-kappa}. Nodes with
$\kappa(v)=0$ (e.g.\ a single child and no observations) do not need a prior
factor.
\end{proof}

For two levels, $\kappa(g)=1-n$ and $\kappa(\ell_j)=0$. For three levels,
$\kappa(g)=1-K$, $\kappa(h_k)=1-n_k$ and $\kappa(\ell_{kj})=0$.

\subsection{The recursive rule}

\begin{theorem}[$L$-level / tree composition rule]\label{thm:tree-rule}
Process the nodes in post-order (children before parents). For each $v$, all
quantities live on the path $\pi(v)$:
\begin{align}
Q_v&=\sum_{o\in O(v)}P_o+\kappa(v)\,P_{0,v}+\sum_{w\in\ch(v)}S_w,\label{eq:tree-Q}\\
\sigma_v&=Q_v^{-1}\Big[\sum_{o\in O(v)}P_os_o+\kappa(v)\,P_{0,v}s_{0,v}+\sum_{w\in\ch(v)}S_w\,\sigma_w^{\alpha(w)}\Big],\label{eq:tree-sigma}\\
S_v&=[Q_v]_{\alpha\alpha}-[Q_v]_{\alpha v}[Q_v]_{vv}^{-1}[Q_v]_{v\alpha},\qquad
K_v=-[Q_v]_{vv}^{-1}[Q_v]_{v\alpha}\qquad(\alpha=\alpha(v)),\label{eq:tree-S}
\end{align}
where $(P_o,s_o)$ are the network's precision and score for observation $o$ on
$\pi(v)$, and $(P_{0,v},s_{0,v})$ are the prior backward precision and diffused
prior score on $\pi(v)$. Then process the nodes top-down:
\begin{equation}\label{eq:tree-down}
u^\rho=\sigma_\rho,\qquad u^v=\sigma_v^v+K_v\big(u^{\alpha(v)}-\sigma_v^{\alpha(v)}\big).
\end{equation}
If every $[Q_v]_{vv}\succ0$ and $Q_\rho\succ0$, this $u$ is the unique solution of
$\Lambda u=r$ with the counting numbers \eqref{eq:tree-kappa}. In particular, it
is the exact diffused tall-data score of every linear-Gaussian tree.
\end{theorem}
\begin{proof}
By induction over the post-order, with the invariant: \emph{when $v$ is
processed, all of $\sub(v)\setminus\{v\}$ has been eliminated, and the factors of
the reduced system that involve block $v$ are exactly the terms in
\eqref{eq:tree-Q}, all with scope $\pi(v)$.} Any original factor that contains
$v$ has scope $\pi(u)$ with $u\in\sub(v)$. Those with $u\neq v$ have been
absorbed by the induction hypothesis into the messages $(S_w,\sigma_w^{\alpha(w)})$
of the children $w$, whose scope is $\alpha(w)=\pi(v)$. Those with $u=v$ are the
observation factors $O(v)$ and the prior factor of $v$. Merging them
(\cref{lem:marginalise}) gives $(Q_v,Q_v\sigma_v)$, the only factor containing
$v$. Marginalising $v$ (\cref{lem:marginalise}, $c=v$, $a=\alpha(v)$) produces the
message $(S_v,\sigma_v^{\alpha(v)})$ on $\pi(\pa(v))$ and the back-substitution
\eqref{eq:tree-down} (an instance of \eqref{eq:backsub}), which maintains the invariant for $\pa(v)$. At the root,
$Q_\rho u^\rho=Q_\rho\sigma_\rho$. Uniqueness: the elimination is a block
$LDL^\T$ factorisation of $\Lambda$ with pivots $[Q_v]_{vv}$ and $Q_\rho$, so
$\Lambda\succ0$ if and only if all pivots are positive definite.
\end{proof}

\paragraph{Interpretation: a subtree is a big observation.}
By \eqref{eq:subtree-sum}, the factors inside $\sub(w)$ cover $\pi(\pa(w))$
exactly once as well. They therefore form the exact tall posterior of the
\emph{sub-model} consisting of the path to $w$, the subtree $\sub(w)$, and its
observations. The message $(S_w,\sigma_w^{\alpha(w)})$ is the Gaussian summary of
that sub-model's composed backward kernel on the parent path: $S_w$ is its
marginal precision on $\pi(\pa(w))$, and $\sigma_w^{\alpha(w)}$ is its composed,
conditional score there. The parent combines its children exactly as the
two-level rule combines observations. This is Gaussian belief propagation, in
precision--score parameterisation, on the tree, run upwards once. The downward
pass is back-substitution, not a second message pass, because only the mean
(Tweedie) is needed.

\paragraph{Cost.} $\sum_v\big(|O(v)|+|\ch(v)|\big)\,d_{\pi(v)}^2+d_{\pi(v)}^3$:
linear in the number of nodes and observations, and cubic in the path dimension
$d_{\pi(v)}\le\sum_{\text{levels}}d_\text{level}$. $\Lambda$ has zero fill-in
because its sparsity graph (the union of root-path cliques) is chordal, and the
leaves-first order is a perfect elimination ordering.

\subsection{Positive definiteness and the information form}
Using $\kappa(v)=1-|O(v)|-|\ch(v)|$,
\begin{equation}\label{eq:tree-info}
Q_v=P_{0,v}+\sum_{o\in O(v)}\big(P_o-P_{0,v}\big)+\sum_{w\in\ch(v)}\big(S_w-P_{0,v}\big)
=:P_{0,v}+\sum_{o}I_o+\sum_wI_w .
\end{equation}
This is the paper's information form, $\Lambda_S=P_0+\sum_jI_j$, at every
level. If every increment is positive semidefinite, then $Q_v\succeq P_{0,v}\succ0$,
all pivots are positive definite, and the rule is well posed. The increments are
PSD in the exact linear-Gaussian case. $I_o$ is PSD because $P_o$ is the precision
of $p(\theta_{\pi(v),0}\mid y_{\pi(v)},x_o)$ and $P_{0,v}$ that of
$p(\theta_{\pi(v),0}\mid y_{\pi(v)})$. $I_w$ is PSD because $S_w$ is the
$\pi(v)$-marginal precision of the sub-model kernel, which additionally
conditions on $y_{\sub(w)}$ and $x_{\sub(w)}$. Gaussian conditioning never
increases covariance. The current PSD projection of $I_j$ therefore extends
directly: project every $I_o$ and $I_w$ at every level.

\subsection{What each method has to supply}
\begin{itemize}
\item \textbf{Observation factors.} The network must output the joint score on
 the whole path $\pi(v)$ for observation $o$. \arrowm{} needs a pilot covariance
 $\Sigma_{0,o}$ of $p(\theta_{\pi(v)}\mid x_o)$ (from the network's own draws), and
 sets $P_o=\Sigma_{0,o}^{-1}+\lam^{-2}\I$. \quiver{} takes
 $P_o=\big[\lam^2(\I+\lam^2\nabla s_o)\big]^{-1}$ from $d_{\pi(v)}$ JVPs or
 finite-difference copies. Both are unchanged; only the dimension grows.
\item \textbf{Prior factors.} For every node with $\kappa(v)\neq0$, the diffused
 prior score $s_{0,v}$ and backward precision $P_{0,v}$ on the \emph{whole path}
 $\pi(v)$ are needed. With two levels, only $\pi(\rho)=\{g\}$ is needed, and
 a Gaussian $p(g)$ makes it closed-form. From three levels on, $p(g,h_k)$
 appears. In partial-pooling models $g$ contains the scale of $h_k$, so this
 prior is a funnel, not a Gaussian, and has no closed-form diffused score. The
 options are: (a) a prior-score network trained on root-path prior draws
 (cheap: no simulator calls), or the same network evaluated with all data
 masked if its training masks include the ``no data'' pattern; (b) its Tweedie
 Jacobian for $P_{0,v}$, as in \quiver; (c) a Gaussian moment surrogate from
 prior samples. Option (c) introduces an error that the two-level rule does not
 have when $p(g)$ is Gaussian. The prior factors carry weights $1-n_k$, so their
 errors are amplified by the group size, exactly as the top-level prior error
 is amplified by $n-1$ today.
\end{itemize}

%=====================================================================
\section{General DAG hierarchies}\label{sec:dag}
%=====================================================================

\Cref{thm:factorisation,prop:linear-system} already cover every model satisfying
\cref{ass:dag}. Exactness of the factorisation therefore holds for every DAG
hierarchy, with no running-intersection or tree assumption, because the regions
are ancestral. What changes is the cost of solving $\Lambda u=r$.

\begin{example}[Crossed random effects]\label{ex:crossed}
$g\to a_i$ ($i\le I$), $g\to b_j$ ($j\le J$), $x_{ij}\mid g,a_i,b_j$, one
observation per cell. The regions are $\{g,a_i,b_j\}$, with intersections
$\{g,a_i\}$, $\{g,b_j\}$ and $\{g\}$. The Möbius numbers are
$c(\{g,a_i,b_j\})=1$, $c(\{g,a_i\})=1-J$, $c(\{g,b_j\})=1-I$, and
\[
c(\{g\})=1-\big[IJ+I(1-J)+J(1-I)\big]=(I-1)(J-1)\;\ge0,
\]
so
\[
p(\theta\mid\X)\propto p(g)^{(I-1)(J-1)}\prod_ip(g,a_i)^{1-J}\prod_jp(g,b_j)^{1-I}\prod_{i,j}p(g,a_i,b_j\mid x_{ij}).
\]
The global prior exponent is \emph{positive}: the row and column corrections
remove $p(g)$ too often, and it has to be multiplied back.
\end{example}

\begin{example}[Shared child]\label{ex:diamond}
$g\to h_1,h_2$; $z_1\mid h_1,h_2$; $z_2\mid h_2$; observations
$x_a,x_b\mid z_1$, $x_c\mid z_2$, $x_d\mid h_1$. The regions are
$\{g,h_1,h_2,z_1\}$ (twice), $\{g,h_2,z_2\}$ and $\{g,h_1\}$. The prior factors
are $p(g,h_1,h_2,z_1)^{-1}$, $p(g,h_1)^{-1}$ and $p(g,h_2)^{-1}$, and
$\kappa(\{g\})=0$.
\end{example}

\subsection{Solving the system}

\begin{enumerate}
\item \textbf{Junction-tree case.} If the maximal regions can be arranged in a
 tree with the running-intersection property, i.e.\ the ancestral observation
 graph is chordal with this clique tree, \cref{thm:tree-rule} carries over with
 nodes replaced by cliques and blocks by clique-minus-separator sets. Each clique
 merges its factors, eliminates its private variables by
 \cref{lem:marginalise}, and passes a Schur message on the separator. The cost is
 linear in the number of cliques and cubic in the clique dimension. Rooted trees
 are the special case where the cliques are root paths.
\item \textbf{General case.} Pick an elimination order (e.g.\ minimum degree)
 and run sparse block Cholesky. Every elimination step is still
 \cref{lem:marginalise}, so only network conditional scores ever enter, but fill-in
 appears. The cost is governed by the treewidth of the ancestral observation
 graph. In \cref{ex:crossed}, after eliminating the $a_i$, the $b_j$ form a dense
 $J\times J$ block: $O(IJ^2+J^3)$, fine for hundreds of levels.
\item \textbf{Large loopy designs.} Under the information-form condition
 $\Lambda\succ0$, conjugate gradients needs only products
 $\Lambda v=\sum_fw_fE_f^\T P_f(E_fv)$, at cost $O(\sum_fd_{R_f}^2)$ per
 iteration. Alternatively, Gaussian BP on $\Lambda$ returns the \emph{exact} mean
 whenever it converges (Weiss \& Freeman, 2001), which is all Tweedie needs.
 Convergence is guaranteed for walk-summable $\Lambda$.
\end{enumerate}

\subsection{Limits of the graph generalisation}\label{sec:dag-limits}

\begin{enumerate}
\item \textbf{Bounded ancestral regions.} The training unit is $(\theta_{R_o},x_o)$.
 If $|R_o|$ grows with the data, as in a latent Markov chain
 $z_1\to z_2\to\cdots$ with $x_t\mid z_t$ where $R_t=\{z_1,\dots,z_t\}$, a
 single-observation network cannot be trained. Two exact alternatives exist:
 \begin{itemize}
 \item[(a)] \emph{Decomposable prior over non-ancestral regions.} If
  $p(\theta)=\prod_Cp(\theta_C)/\prod_Sp(\theta_S)$ over a junction tree
  (e.g.\ $C_t=\{z_{t-1},z_t\}$, $S_t=\{z_t\}$), and the network for $x_o$ is
  trained on $(\theta_{C(o)},x_o)$ with $\theta_{C(o)}$ drawn from the
  \emph{true prior marginal} $p(\theta_{C(o)})$ (e.g.\ cropped windows of long
  simulations), the proof of \cref{thm:factorisation} goes through with these
  regions. The counting condition is the junction-tree identity.
 \item[(b)] \emph{Training-prior ratio form.} For any regions $R_o\supseteq\pa(o)$
  and any training prior $\tilde p_o(\theta_{R_o})$,
  $p(\theta\mid\X)\propto p(\theta)\prod_o\tilde p_o(\theta_{R_o}\mid x_o)/\tilde p_o(\theta_{R_o})$.
  Each coordinate is covered $1+\sum_o(1-1)=1$ times, so
  \cref{thm:factorisation,prop:linear-system} hold with
  $\Lambda=P_0^{\text{full}}+\sum_oE_o^\T(P_o-\tilde P_{0,o})E_o$. This needs the
  diffused prior score and precision on all of $\theta$, which are tractable for
  Gaussian Markov priors (sparse $P_0^{\text{full}}$) but not in general.
 \end{itemize}
\item \textbf{Observations must be sinks.} Models in which data feed back into
 latents or other data (autoregressive observations) are outside
 \cref{ass:dag}.
\item \textbf{The single approximation is unchanged, but it is applied in
 higher dimension.} Each observation factor is a Gaussian surrogate on
 $d_{R_o}$ dimensions instead of $d_g+d_\ell$. The multimodality and
 heavy-tail failure modes documented for the two-level rule (e.g.\ the IK branch
 bias summed $n$ times) apply per factor. Prior factors appear at every
 intermediate level with weights $1-|O(v)|-|\ch(v)|$, so prior-score and
 Gaussian-surrogate errors are amplified at each level, not only at the root.
\end{enumerate}

%=====================================================================
\section{Numerical verification}\label{sec:numerics}
%=====================================================================

A standalone NumPy script (CPU only, 3 threads; not part of the package) builds
random linear-Gaussian DAGs, samples data, and computes the exact Gaussian
posterior of every factor under the true prior marginal of its region, together
with the exact diffused tall-data score $-(C+\lam^2\I)^{-1}(y-m)$. It then
compares (i) the dense solve $\Lambda^{-1}r$ with Möbius counting numbers,
(ii) the recursive tree rule of \cref{thm:tree-rule} (trees only), and (iii) an
ablation that keeps only the root prior correction and drops all intermediate
prior factors. It also asserts \eqref{eq:cover} for every node. The table
reports relative error $\|u-u_\text{exact}\|/\|u_\text{exact}\|$, maximised over
$\lam\in\{0.01,0.3,1,5\}$ with a fresh random $y$ per $\lam$.

\begin{center}\small
\setlength{\tabcolsep}{4pt}
\begin{tabular}{@{}p{3.9cm}rrccc p{2.9cm}@{}}
\toprule
model & dim & obs & dense & tree rule & drop mid priors & non-zero $\kappa$\\
\midrule
3-level tree ($K=3$, $n_k=2,3,4$, obs on $h_0$, 2 trials on $\ell_{21}$) & 23 & 11 & $2.2\times10^{-15}$ & $2.3\times10^{-15}$ & 4.8\% -- 48\% & $g{:}{-}2$, $h_k{:}{-}2,{-}2,{-}3$, $\ell_{21}{:}{-}1$\\
5-level random tree & 46 & 13 & $6.7\times10^{-15}$ & $6.7\times10^{-15}$ & 3.3\% -- 21\% & 11 path priors\\
crossed $3\times4$ (\cref{ex:crossed}) & 9 & 12 & $2.4\times10^{-15}$ & -- & 8.6\% -- 67\% & $g{:}{+}6$, $a_i{:}{-}3$, $b_j{:}{-}2$\\
shared child (\cref{ex:diamond}) & 5 & 4 & $4.7\times10^{-16}$ & -- & 0.8\% -- 12\% & three ${-}1$\\
\bottomrule
\end{tabular}
\end{center}

The derived counting numbers are therefore exact on all four structures, and the
recursive tree rule agrees with the dense solve to round-off. The ablation error
is largest at small $\lam$: the intermediate prior corrections are not optional,
and are not a high-noise effect. The crossed design also confirms the positive
global exponent $\kappa(g)=+6=(3-1)(4-1)$.

%=====================================================================
\section{Summary}\label{sec:summary}
%=====================================================================

\begin{enumerate}
\item \textbf{Three levels: yes.} \Cref{eq:pass1,eq:pass2-merge,eq:pass2,eq:pass3}.
 The new required inputs are the network's path score and precision on
 $(g,h_k,\ell_{kj})$, and the diffused prior score and precision on $(g,h_k)$,
 which enter with weight $1-n_k$.
\item \textbf{$L$ levels / any rooted tree: yes.} \Cref{thm:tree-rule} is the
 two-level rule applied recursively, with subtrees acting as observations. It has
 one post-order pass and one back-substitution, cost linear in the number of
 nodes, and an information form with a per-level PSD safeguard. It is exact on
 linear-Gaussian trees.
\item \textbf{Any DAG hierarchy with sink observations: yes for exactness, with
 structure-dependent cost.} The counting-number factorisation
 (\cref{thm:factorisation}) is exact for every DAG because the regions are
 ancestral. The composed score is always $\Lambda^{-1}r$ built from network
 conditional scores. Junction-tree structures keep a closed-form message rule;
 loopy designs need a sparse or iterative solve.
\item \textbf{Main practical obstacles:} (i) diffused prior scores on intermediate
 paths (non-Gaussian funnels in partial pooling); (ii) the network input
 dimension grows with depth; (iii) unbounded ancestral sets (state-space models)
 need the non-ancestral or ratio variants of \cref{sec:dag-limits}.
\end{enumerate}

\end{document}
